%***************************************************************************************************
% preamble/packages.tex - lista pakietów w ustalonej kolejności.
% Kolejność ładowania jest istotna - nie zmieniaj jej bez potrzeby.
% Pakiet hyperref jest ładowany osobno, na samym końcu (preamble/links.tex).
%***************************************************************************************************

% ---------------------------------------------------------------------------------------------
% Język, kodowanie, czcionki, typografia (dokument wyłącznie w języku polskim)
% ---------------------------------------------------------------------------------------------
\usepackage[T1]{fontenc}
\usepackage[utf8]{inputenc}    % wymagane przez klucz inputencoding pakietu listings
\usepackage{lmodern}
\usepackage[polish]{babel}
\let\lll\undefined             % babel (polish) definiuje \lll, co koliduje z amssymb
\usepackage{csquotes}          % polskie cudzysłowy: \enquote{tekst}
\usepackage{indentfirst}       % wcięcie pierwszego akapitu po nagłówku (zwyczaj polski)
\usepackage{microtype}
\frenchspacing                 % polska typografia: jednakowe odstępy po kropce

% ---------------------------------------------------------------------------------------------
% Matematyka
% ---------------------------------------------------------------------------------------------
\usepackage{amsmath}
\usepackage{amssymb}
\usepackage{amsthm}
\usepackage{mathtools}

% ---------------------------------------------------------------------------------------------
% Kolory, grafika, obiekty pływające
% ---------------------------------------------------------------------------------------------
\usepackage[table, svgnames]{xcolor}
\usepackage{graphicx}
\usepackage{float}             % umożliwia [H] (dokładne umiejscowienie obiektu)
\usepackage{caption}
\usepackage{subcaption}        % podrysunki: środowisko subfigure

% ---------------------------------------------------------------------------------------------
% Tabele
% ---------------------------------------------------------------------------------------------
\usepackage{array}
\usepackage{tabularx}
\usepackage{booktabs}
\usepackage{longtable}
\usepackage{multirow}
\usepackage{makecell}
\usepackage{seqsplit}          % łamanie długich ciągów (np. skrótów SHA-256) w dowolnym miejscu

% ---------------------------------------------------------------------------------------------
% Listy, listingi, algorytmy, dołączanie PDF
% ---------------------------------------------------------------------------------------------
\usepackage{enumitem}
\usepackage{listings}
\usepackage[ruled, vlined, linesnumbered, longend, algochapter]{algorithm2e}
\usepackage{pdfpages}          % \includepdf - wstawianie stron z plików PDF (np. do załączników)
\usepackage{multicol}
\usepackage{pifont}            % znaczniki \ding (\cmark, \xmark w makrach)

% ---------------------------------------------------------------------------------------------
% Układ strony, nagłówki i stopki, tytuły rozdziałów
% ---------------------------------------------------------------------------------------------
% Marginesy 25 mm; bindingoffset zostawia dodatkowe 5 mm na zszycie (przy druku dwustronnym).
\usepackage[a4paper, top=25mm, bottom=25mm, left=25mm, right=25mm,
            bindingoffset=5mm, headheight=30pt, headsep=12pt]{geometry}
\usepackage{fancyhdr}
\usepackage{titlesec}
